This synthesis companion omits the complete appointed-text component and the
ten-part documentary sweep. It offers instead a sustained reading of every
proper element from the common research record.

\subsection*{Gift creates the servant who can run \enspace\properrefs{Int., Coll., Ep., Postcomm.}}

The Mass begins before human sufficiency. David's urgent
\latin{Deus, in adiutórium} gives the worshipper the speech of dependence;
Eusebius and Augustine hear in it Christ and his members, while Cassian and
Benedict make it the prayer of every condition. The Collect then discloses
what that help does. Worthy service comes \latin{de cuius múnere}, yet the
faithful really serve; God gives, yet they really run
\latin{ad promissiónes tuas}. Grace does not occupy the space where creaturely
action should have been. It heals and elevates the creature so that filial
action becomes possible.

Paul names the same economy ministerially. Confidence travels through Christ
toward God, sufficiency is not \latin{quasi ex nobis}, and nevertheless God
has made real ministers of the new covenant. Chrysostom points to the
Corinthians as a Spirit-written apostolic letter; Theodoret and Aquinas insist
that God's causality establishes rather than empties ministerial agency.
Augustine's exposition of Psalm 118 joins the enlarged-heart run, love of God
and neighbor, letter and Spirit, and continual praise. The Sunday therefore
possesses a coherent anthropology: the cry receives help; help gives service;
service becomes a course; the course becomes ministry.

The Postcommunion completes that line by returning to gift at a deeper level.
Holy \latin{participátio} is the subject of \latin{vivíficet} and
\latin{tríbuat}; the mystery gives life, expiation, and defense. Ignatius calls
the Eucharist medicine of immortality, and Aquinas unfolds its nourishing,
sin-remitting, unitive, and protective fruits. The faithful are not spectators:
they cry, assent, serve, run, minister, offer, receive, and go. Nor are they
self-saviors: God attends, makes haste, grants, makes sufficient, pardons,
vivifies, and guards. Catholic cooperation is the motion of a living member in
Christ.

\subsection*{The written Law reaches life in Christ the Neighbor \enspace\properrefs{Ep., Gosp., Coll.}}

Paul's ``letter kills'' and Christ's ``What is written?'' interpret one another.
The old ministry is divine and glorious; the command to love God and neighbor
is true. What external inscription cannot do by itself is forgive the sinner
and write charity within. Ambrosiaster locates the new ministry's glory in
forgiveness and righteousness; Augustine teaches that even the good command is
killing letter without the inward Spirit; Aquinas therefore calls the grace of
the Holy Spirit the New Law's principal reality. The Gospel does not evade the
letter but displays its Spirit-given perfection.

The lawyer's joined citation of Deut.~6:5 and Lev.~19:18 is correct. His effort
to justify himself appears when he asks which persons fall inside the
neighbor-boundary. Christ replies with deeds: compassion, approach, binding,
oil and wine, bearing, shelter, vigil, payment, entrusted care, and return. The
question is reversed from identifying a neighbor to becoming one. Hence
\latin{Hoc fac, et vives} is neither autonomous moralism nor a dismissal of
doctrine. It is the good command reaching life in a heart acted upon by grace.

The received allegory answers whence such life comes. Origen transmits the
elder's vision of wounded humanity; Ambrose, Augustine, and Bede contemplate
the Samaritan as Christ. Humanity descends from Jerusalem into mutable
mortality, loses the robe of grace, is wounded by sin, and lies half-alive.
The Son becomes near in the Incarnation, binds the wounds, bears weakness in
his flesh, brings the patient to the Church, entrusts care to apostolic
ministry, and promises to return. Oil becomes pardon and consolation; wine,
compunction and fervent exhortation; the morrow opens to Resurrection; the
return to Parousia and judgment.

Bede states the unity of senses: every doer of mercy becomes neighbor, and
Christ is Neighbor above all. Those healed by Christ are not excused from
\latin{fac simíliter}; his charity becomes active in them. The Collect's run
and the Gospel's go are therefore one motion. Grace approaches fallen humanity,
puts it upon the road, and makes the recipient capable of approaching another.
The end is the eternal life asked about at the beginning: the Church as inn is
a pilgrim household awaiting the Lord's return and the heavenly patria.

\subsection*{Four psalm chants consecrate peril, praise, night, and provision \enspace\properrefs{Int., Grad., All., Comm.}}

The four chants teach the Christian mouth across the whole span of creaturely
life. The Introit speaks from peril and asks haste. The Gradual blesses in
every time and glories only in the Lord. The Alleluia prays before the God of
salvation by day and night. The Communion praises the Creator whose works fill
earth and yield bread, wine, and oil. Emergency, deliverance, darkness, and
abundance become forms of one Godward address.

Patristic reception deepens each register in Christ. Eusebius hears David,
Christ, and the Church in Psalm 69; Augustine hears the whole Christ confessing
perpetual need. Basil makes Psalm 33's ``always'' embrace prosperity and loss.
Augustine reads the psalm's Abimelech title beside the Old-Latin/Septuagint-shaped
1~Samuel 21:13 lemma ``he was carried in his own hands'' and sees Christ
carrying the sacrament of his own Body; the clause belongs to Augustine's
Samuel text, not to the psalm title or the appointed Gradual. In Psalm 87
Augustine hears the Passion and the Church praying in
prosperity and adversity; Eusebius hears the incarnate Lord's continuous
oblation; Theodoret lets afflicted Israel and suffering humanity pray. The
Alleluia is therefore not praise pasted over lament: it is Paschal praise from
within the Lord's night.

Psalm 103 then brings the entire material order into thanksgiving. Theodoret
rejoices in the bodily service of bread, wine, and oil. Eusebius passes from
created goods to the true Vine, mystical Logos, illumination, and healing.
Augustine reads apostolic earth bearing Christ as bread, the cup's sober
inebriation, and the anointing of the whole Christ; Cassiodorus explicitly
names wine consecrated as Christ's Blood and oil as royal chrism. At Communion
the created gifts are not erased: they reach their highest service when the
Church offers what earth and human labor provide and receives Christ himself
under sacramental signs.

The chant sequence thus gives a spirituality capacious enough for the Gospel.
The Christian can cry without despair, praise without denying affliction, pray
through night without ceasing to name God salvation, and receive abundance
without forgetting the Creator or the wounded neighbor. The meek rejoice in a
boast directed to God; the communicant receives a gift ordered to thanksgiving
and generous service.

\subsection*{Intercession and compassion meet in Christ's one oblation \enspace\properrefs{Gosp., Off., Sec.}}

The Gospel and Offertory place two traditional figures of Christ on different
axes of mercy. The Samaritan moves toward wounded humanity: he sees, approaches,
binds, bears, pays, and promises to return. Moses stands before God for a
guilty people: he invokes deliverance from Egypt, the divine name among the
nations, and the oath to Abraham, Isaac, and Jacob. One movement is toward the
sufferer; the other is intercession before judgment. Both are fulfilled in the
Mediator who assumes our condition, gives himself for sinners, heals, and ever
lives to intercede.

The Exodus tradition reads Moses with full theological weight. Augustine
explains that threatened evil is penal judgment, not moral evil in God, and
that providence includes the intercessor's real prayer. Gregory hears ``let me
alone'' as God granting Moses boldness and Moses' willingness to perish as a
self-offering. Tertullian explicitly calls him a type of Christ. Theodoret
preserves both mercy and judgment: pardon protects the covenantal future while
discipline remains. Mercy is not a revision forced upon God; the immutable God
raises the mediator through whom his faithful purpose is accomplished.

The Secret carries this typology to the altar. God first regards the offerings
\latin{propítius}; then those offerings give honor by granting pardon.
Justin, Irenaeus, Cyril of Jerusalem, and Augustine place created gifts,
thanksgiving, ecclesial self-offering, and intercession within Christ's
sacrifice. Trent supplies the decisive confession: the Mass is truly
propitiatory because the same Christ who offered himself once on Calvary is
offered sacramentally in an unbloody manner. The offerings neither buy mercy
nor merely illustrate it. Taken into Christ's one action, they are the true
sacrificial offering through which the Church asks and receives pardon.

This is why sacrifice and neighbor-love cannot be separated. Christ's
oblation forms the Church as his merciful body. Forgiven worshippers honor the
divine name by praise and by becoming neighbors; works of mercy, ordered to
God, become spiritual sacrifices of Christ's members. The lawyer's
self-justification yields to pardon; pardon yields to gratitude; gratitude
takes bodily form in the care of the wounded. At the altar the Church learns
both directions of the Mediator's love---toward the Father in obedient
self-offering, and toward fallen humanity in healing compassion.

\subsection*{The half-dead is made alive and fortified \enspace\properrefs{Ep., Gosp., Comm., Postcomm.}}

The formulary's central verbal line is life. Robbers leave the traveller
\latin{semivívo}; Christ says \latin{Hoc fac, et vives}; Paul proclaims
\latin{spíritus autem vivíficat}; after Communion the Church asks
\latin{Vivíficet nos \ldots{} huius participátio sancta mystérii}. Correct
reading is necessary but cannot heal fallen nature by itself. The Spirit gives
life, Christ draws near, the Church shelters, and sacramental participation
nourishes the restored person.

The Postcommunion does more than repeat ``life.'' Its one participation gives
\latin{expiatiónem} and \latin{munímen}: it cleanses the wound of sin and arms
the recipient against renewed assault. The Mass began with a cry amid enemies
and ends by asking for defense. It began with God hastening to help and moved
through the faithful running, the Samaritan approaching, Moses interceding, and
the gifts being offered. Communion is the sacramental center from which these
motions receive endurance.

Aquinas provides the theological map in \work{ST} III, q.~79, aa.~1--2, 4, 6,
8: the Eucharist gives and nourishes the life of grace, remits venial sin,
strengthens charity, guards against future sin, unites the members to Christ
and one another, and pledges glory. The gift is objective and Christ's; its
fruit flowers in worthy reception. Hence the last petition is both medicinal and missionary. Those once half-dead are made
alive in Christ; those borne by him become bearers of burdens; those cleansed
seek pardon for sinners; those fortified resist the fear and self-protection
that pass by on the other side.

The promised end gathers the entire Mass. The Collect names God's promises as
the runner's destination. Moses pleads the patriarchal promise. Christ the
Samaritan promises to return. The Postcommunion gives the life and defense by
which pilgrims continue until that return. Thus the proper is neither a moral
lesson appended to sacrament nor a sacramental consolation detached from moral
life. It is one movement of mercy: the Father gives, the Son approaches and
offers, the Spirit vivifies, the Church receives and ministers, and healed
persons go forth to love on the road toward eternal life.
